\chapter{The Maintenance Playbook}
\label{ch:playbook}

A cookbook, organised as "if you need to do X, start here." Each entry names the exact files to touch and the chapter with the full background.

\section{"I need to add a new page"}
\begin{enumerate}
  \item Create \pathname{src/app/(dashboard)/your-route/page.tsx} (or outside the group if it shouldn't share the dashboard shell — Chapter~\ref{ch:proxy}).
  \item Decide Server vs. Client Component \emph{first}, based on whether it needs interactivity (Chapter~\ref{ch:architecture}) — don't default to \code{"use client"} out of habit; most new read-only pages should be Server Components even though most existing ones aren't (Chapter~\ref{ch:pages}).
  \item If it should be reachable only when logged in, add its path prefix to \code{PROTECTED\_ROUTES} in \pathname{src/proxy.ts} (Chapter~\ref{ch:proxy}) — this is a manually maintained array, nothing enforces the addition for you.
  \item Add a link in \pathname{dashboard-sidebar.tsx} if it belongs in primary nav (Chapter~\ref{ch:components}).
\end{enumerate}

\section{"I need to change the database schema"}
\begin{enumerate}
  \item \code{npx supabase migration new descriptive\_name} (Chapter~\ref{ch:workflow} — never hand-create the file).
  \item Write the migration with a comment block explaining \emph{why}, following the house style in \pathname{20260701182526\_...} and \pathname{20260701182527\_...} (Chapter~\ref{ch:workflow}).
  \item \code{npx supabase db reset} locally, and manually re-exercise every page that touches the changed table (Chapter~\ref{ch:workflow} — the test suite is minimal and does not cover page-level flows yet).
  \item If the change touches \code{books} or \code{clubs}, check whether \code{get\_books\_nearby}/\code{get\_clubs\_nearby} need their \code{RETURNS TABLE} column list updated too (Chapter~\ref{ch:functions} — this has bitten this codebase once already). If instead you're changing a function's \emph{parameter} list, \code{DROP FUNCTION} the old signature explicitly before \code{CREATE OR REPLACE} — a new parameter list creates a separate overload rather than replacing the function, which has also bitten this codebase once already (Chapter~\ref{ch:functions}'s second \code{get\_books\_nearby} gotcha). Verify with \code{\textbackslash df public.your\_function} that exactly one overload exists afterward.
  \item Update Chapter~\ref{ch:schema} of \emph{this manual} in the same pull request.
\end{enumerate}

\section{"I need to add a new admin action"}
Add a function to \pathname{src/lib/admin-actions.ts} following the guard/act/log/revalidate shape exactly (Chapter~\ref{ch:server-actions}): \code{guardRole(minRole)} first, the mutation(s), an \code{auditLog(...)} call, \code{revalidatePath('/admin')}, and a \code{try/catch} returning \code{\{success, error?\}}. Pick \code{minRole} deliberately — \code{'moderator'} for anything reversible, \code{'super\_admin'} for anything destructive or that changes another admin's privileges (Chapter~\ref{ch:rls}).

\section{"I need to add a new notification type"}
\begin{enumerate}
  \item Call \code{createNotification} (Chapter~\ref{ch:server-actions}) with a new \code{type} string from wherever the triggering event happens.
  \item If it should also send an email, add an entry to \code{NOTIFICATION\_SUBJECTS} in \pathname{send-notification-email.ts} (Chapter~\ref{ch:server-actions}) — falls back to a generic subject if you skip this, so it degrades gracefully, but do it anyway.
  \item Add an icon/label for it in \emph{both} places that maintain that mapping independently: \pathname{notifications/page.tsx} and \pathname{notification-bell.tsx} (Chapter~\ref{ch:pages} flags these as already-drifted duplicates — adding a third divergence makes this worse, so consider extracting a shared constant while you're in there).
  \item Ensure any user-influenced text passed into the notification's \code{title} goes through \code{escapeHtml} before it can reach the email HTML template (Chapter~\ref{ch:server-actions}).
\end{enumerate}

\section{"I need to change a rate limit or feature flag"}
Edit the row in \code{platform\_settings} — either through \pathname{admin/settings/page.tsx} directly, or a migration's \code{INSERT ... ON CONFLICT (key) DO UPDATE} if it should ship as a new default (Chapter~\ref{ch:functions}). Remember \code{get\_platform\_setting\_int} expects the stored JSON to be either a JSON string or JSON number — writing raw unquoted text into the \code{value} column outside the app's own read/write path can break the \code{\#>> '\{\}'} unwrap (Chapter~\ref{ch:functions}).

\section{"I need to test something as a banned user / admin / moderator"}
Register a fresh account locally (Chapter~\ref{ch:local-env}), then flip the relevant \code{profiles} columns directly in Studio's table editor or via SQL — \code{is\_banned = true}, or \code{is\_admin = true, admin\_role = 'moderator'} (both fields, together — Chapter~\ref{ch:rls} explains why one alone grants nothing). This is faster and more reliable than trying to reach the same state through the UI.

\section{"I need to regenerate the schema dump"}
\code{npx supabase db dump --local --schema public > supabase/schema.sql} (Chapter~\ref{ch:workflow}), after your local database reflects every migration you want captured. Do this deliberately, as its own commit, not as an incidental side effect of an unrelated change.

\section{"I need to run a security \& code-quality audit"}
This isn't a one-time event described elsewhere in this manual as history — it's a practice worth repeating periodically (monthly is a reasonable cadence for a project this size), and the process is the same each time:

\begin{enumerate}
  \item \textbf{Survey before fixing.} Look across at least three angles, since they catch different things: leaked secrets/credentials (grep the tree and git history, check \pathname{.gitignore} covers every \pathname{.env*} file, check client-vs-server secret boundaries per Chapter~\ref{ch:server-actions}); code quality (god files, duplicated patterns, dead code — Chapter~\ref{ch:workflow}'s lint-debt cleanup is exactly this category); and RLS/database security, which means reading every migration \emph{in order} rather than any single one in isolation (Chapter~\ref{ch:rls}) — a policy can look correct in the migration that added it and still be wrong in the schema's current, cumulative state, the way \pathname{books}' legacy delete policy was.
  \item \textbf{Verify a finding yourself before trusting it, especially a severe one.} Don't just accept a claimed vulnerability — reproduce it. The profile self-escalation issue (Chapter~\ref{ch:rls}) was confirmed by reading the actual \code{GRANT}/policy statements in the migration files before writing a single line of fix, not by taking a description of the bug at face value.
  \item \textbf{Fix in severity order, and get explicit sign-off before anything touches the linked cloud project.} Write migrations locally, run \code{supabase db reset} (Chapter~\ref{ch:workflow}), and prove the fix with SQL that simulates \emph{both} the attack and the legitimate flow it must not break — \code{SET ROLE authenticated; SET request.jwt.claims = '...'} to impersonate a specific test user's session in the local Studio SQL editor (Chapter~\ref{ch:rls}'s intern exercise shows the \code{anon}-role version of this same technique) is how every fix in this pass was actually proven, not assumed. Only run \code{supabase db push} once local proof exists, and only against the local stack's own linked project — never experiment with attack payloads against the hosted database.
  \item \textbf{Verify code-quality fixes by driving the real app, not just by reading the diff.} \code{tsc --noEmit}, \code{npm run lint}, and \code{npm run build} catch a lot, but none of them prove a refactored page still works end to end. Point the dev server at the \emph{local} Supabase stack specifically for this (Chapter~\ref{ch:local-env}) — a headless browser script that logs in as a real local test account and exercises the actual flow (add a book, accept a request, whatever the diff touched) is the only thing that catches a prop wired to the wrong callback after a component split.
  \item \textbf{Update this manual in the same pass}, per the entry below — not as an afterthought once the code is merged.
\end{enumerate}

\section{"Something changed in this codebase and this manual is now wrong"}
Fix the manual in the same pull request as the code change, in the relevant \pathname{Guide/chapters/*.tex} file. This is not a nice-to-have: Chapter~\ref{ch:schema}'s note about \pathname{schema.sql} having gone five migrations stale is a direct, cautionary illustration of what happens to any documentation artifact that isn't updated as part of the change it describes. Don't let this manual become the next example.

\begin{olktask}
Pick any one row in this playbook you haven't personally done yet, and do it against your local stack, start to finish, before reading further into this manual. Muscle memory beats reading every time.
\end{olktask}
